\question[4]
Automaten im Alltag.
\begin{parts}
\part Nenne vier Geräte oder Programme, deren Verhalten man mit einem Automatenmodell beschreiben kann.
\part Wähle einen Getränkeautomaten. Beschreibe, was bei ihm die Zustände, die Eingaben und die Ausgaben sind.
\part Erläutere, warum ein Automat für die Suche eines Wortes in einem Text geeignet ist.
\end{parts}

\begin{solutionbox}{6cm}
\begin{parts}
\part z.B. Ampel, Getränke-/Snack-/Ticketautomat, Leergutautomat, Aufzug, Waschmaschine,
Login-Dialog oder Buttons einer grafischen Oberfläche, Netzwerkprotokolle (TCP), Textsuche.
\part Zustände: bisher eingeworfener Betrag (z.B. 0\,ct, 50\,ct, 100\,ct).
Eingaben: eingeworfene Münzen, Tastendruck für die Getränkewahl, Abbruchtaste.
Ausgaben: Getränk, Wechselgeld, Anzeige auf dem Display.
\part Der Automat liest den Text Zeichen für Zeichen. Er muss sich nur merken, wie viel vom Suchwort er
gerade schon gefunden hat. Das sind endlich viele Möglichkeiten, also endlich viele Zustände.
Der Text muss nur einmal von vorne nach hinten gelesen werden.
\end{parts}
\end{solutionbox}
